\section{The Ladder Under Load: What a Star Does When It Loses}

There is a complaint I would make against this book if I were reading it rather than writing it. It has argued, for six chapters, that the world is layered without end --- and it has never once shown the tower under stress. Every picture so far has been of a universe going quietly about its business: rain falling, shadows overlapping, crowds holding their shape. A tower that only has to stand still is not much of a tower. What happens when you lean on it?

The universe runs that experiment for us, on a schedule, in public. The apparatus is called a star, and the experiment begins when the star runs out of fuel.

A star is a long argument between two pressures: the crowd pressing it inward, and the heat of its own burning pressing outward. When the fuel is gone the heat fails, and the star falls inward until something else stops it. Something else does. In a star of modest mass the something is the electrons: crowded past a certain closeness they refuse to be crowded further --- not because they repel one another, though they do, but because the world will not permit two of them into the same state. That refusal is stiff enough to hold the mass of the sun in a ball the size of the Earth.\footnote{Degeneracy pressure: the stiffness that arises because identical particles of certain kinds cannot occupy the same state. It is not a force between them; it is a prohibition on their arrangement. (Dictionary.)} The result is a white dwarf, and our own sun will make one.

Load it further and the electrons lose. Above about one and a half times the mass of the sun --- Chandrasekhar computed the figure at nineteen, on the voyage from India to England --- the electron strut buckles, electrons are crushed into protons, and what stands up out of the wreckage is a body made almost entirely of neutrons: the mass of a sun inside the space of a city. The neutrons hold in their turn, by the same refusal plus the shove of the strong force, until something over two solar masses, where the neutron strut buckles too.

And then, in the standard picture, there is nothing left. No third strut is known. The equations run on to a place where the density is infinite and the volume is nothing --- a singularity\footnote{Singularity: a point in a theory's predictions where a quantity becomes infinite. Physicists generally read such a point as the theory failing rather than as a description of something real.} --- and the physicists who use the word will tell you, without being pressed, that it is not a description of anything. It is the mark left where the theory stopped working.

Look at the shape of that story, because the shape is the whole point. Each rung is the same event at a finer scale: a structure's resistance is exceeded, its constituents are crushed into a smaller arrangement, and the smaller arrangement's own stiffness takes the weight. Heat fails; electrons hold. Electrons fail; neutrons hold. That is not an analogy for this book's thesis. It is this book's thesis, observed, twice, in objects we can point a telescope at.

And now notice what is doing the holding, because this is where I must be careful with my own enthusiasm. Degeneracy pressure is not a force in the ordinary sense. The standard account calls it exclusion: two identical particles of certain kinds are not permitted into the same state, and the not-permitting is stiff enough to carry a star. As physics that is not in dispute --- the white dwarfs are up there. But as an \emph{explanation} it is, in this book's terms, unfinished, because a prohibition is not a push. A rule that forbids, with nothing doing the forbidding, is the same species of object as a force that pulls across empty space with nothing in between; I have spent seven chapters declining that, and I cannot take the decline back here merely because the prohibition is holding up something I want.

So let me write the debt where debts belong, in the open. If this book is right, exclusion is not a decree handed down but a congestion: something about two identical arrangements of the crowd that cannot be pressed into one region without the crowd itself objecting, and objecting harder the closer they are pressed. The materials for such an account are already lying in Chapter~6 --- matter, here, is circulation, and circulations sharing a region are not a neutral matter. I do not have the account. I want to say so in the same breath as the claim, because this ladder is the best evidence the book has, and evidence that rests on an unexplained prohibition is evidence with a hole in the middle of it.

What survives is weaker and still worth having. Whatever the stiffness turns out to be made of, it evidently comes layer by layer: each rung has its own version of it, and therefore its own limit mass. This ladder of collapse is the tower seen edge-on, under load.

\begin{figure}[htbp]
\centering
\begin{tikzpicture}[>=stealth]
  % the load
  \draw[->, bbuBlue, line width=1.6pt] (0.55,7.45) -- (0.55,0.30);
  \node[anchor=south, bbuBlue] at (0.55,7.55) {\small the load};
  % observed rungs
  \foreach \y/\body/\strut in {%
    6.5/a star/heat holds,
    4.6/a white dwarf/electron degeneracy holds,
    2.7/a neutron star/neutron degeneracy holds} {
    \draw[line width=0.9pt] (1.35,\y) rectangle (5.65,\y+0.8);
    \node at (3.5,\y+0.4) {\small \body};
    \node[anchor=west, bbuGold, align=left, text width=2.9cm] at (5.9,\y+0.4) {\small \strut};
  }
  % the unknown rung
  \draw[bbuGold, dashed, line width=0.9pt] (1.35,0.8) rectangle (5.65,1.6);
  \node[align=center, text width=3.9cm] at (3.5,1.2) {\small a body, held --- one floor down, or two};
  \node[anchor=west, bbuGold, align=left, text width=2.9cm] at (5.9,1.2) {\small a strut we cannot name};
  % failures
  \foreach \y/\why in {6.1/fuel gone, 4.2/above $\approx 1.4\,M_\odot$, 2.3/above $\approx 2\,M_\odot$} {
    \draw[->, bbuRed, line width=1.0pt] (3.5,\y+0.30) -- (3.5,\y-0.30);
    \node[anchor=west, bbuRed] at (3.72,\y) {\small \why};
  }
  \node at (3.5,0.42) {$\vdots$};
  \node[anchor=north, align=center, text width=6.8cm] at (3.5,0.15)
    {\small how much further, we cannot see from here};
\end{tikzpicture}
\caption{The ladder under load. Each rung is the same event at a finer scale: a structure's resistance is exceeded, its constituents are crushed into a smaller arrangement, and that arrangement's own stiffness takes the weight. The first three rungs are observed --- our sun will make the second, and we have photographed the third. The standard picture has no fourth rung and puts a singularity there instead: everything, in no volume at all, which is a floor by another name. The tower's fourth rung is an ordinary thing in extraordinary circumstances --- a body, held up by the layer below.}
\label{fig:collapseladder}
\end{figure}

Which makes the top of the ladder the interesting place. The standard picture, past the neutron rung, offers a completed infinity: everything, in no volume at all. I have spent this book declining exactly that kind of object --- not because infinities are distasteful, but because \emph{and here it stops, and do not ask what it is made of} is the move Chapter~1 called cheating when it was an edge, and Chapter~3 called arbitrary when it was a floor. A singularity is a floor. It is the most extreme floor anyone has ever proposed.

The tower does not need it. Past the neutron rung there is another rung: the layer below, taking the load on its own strut. The matter is not destroyed and does not become infinite. It is crushed into an arrangement our layer has no words for, and held up by a stiffness our layer cannot measure.

Here I want to be careful, because there is an overclaim available and I do not want it. That the tower is endless is a claim about the \emph{world}; how far a particular star falls is a claim about an \emph{event}. The two use the same word and are not the same claim, and this book asserts the first. The tower has no floor.

The second rests on the second rule of Chapter~6 --- that what happens in a layer is local to that layer and its neighbours --- which is itself a conjecture and was labelled one where it was stated. If that rule holds, a collapse is not a plunge down an open shaft. It is a sequence of separate events, each one a floor being overwhelmed and handing the load to the floor beneath, and there is no reason for such a sequence to run forever and every reason to expect it to stop at the first floor stiff enough to take the weight. A black hole would then be a body that has fallen a rung, or a few, and is sitting there, held --- exactly as a neutron star sits held one rung below a white dwarf. It has a state. It has structure. It is simply not our layer's structure.

I believe that firmly, and I want to be exact about the standing of it: I cannot prove it, and I am not claiming the alternative impossible. I made that mistake once already, in Chapter~3, where I went in to prove that no bottom was \emph{possible} and had the impossibility taken off me by an opponent who was right to take it. A conjecture held firmly and an impossibility established are different objects, and confusing them is the one error this book has already paid for once.

And from the outside the difference does not show. A body that has descended one floor and a body descending without end present the same face: dark, dense, and censored. The horizon does not report the number. So the modest claim is the one the evidence carries --- some rungs below ours, and we cannot say how many --- and anyone who tells you the number is telling you more than they have. That includes the singularity, which is the claim that the number is infinite.

Which suggests the name is wrong. \emph{Black hole} was coined inside a picture in which there is nothing at the centre but a point of infinite density --- a hole, in the plain sense that there is nothing there. In this book's picture there is a great deal there: a very dense body, made of the layer below, doing what dense bodies do. Not a hole. A star a floor or two down, wearing a horizon.

And the horizon itself gets demoted, by machinery this chapter has already bought and paid for. The second wound gave every layer its own light and its own speed limit; the fifth wound spent that coin, buying channels below ours that outrun our light. Apply it here and the famous sentence changes meaning. \emph{Nothing escapes} means nothing escapes \emph{at our layer's speed} --- and the layer below is not bound by our layer's speed. Its light is faster. A horizon is a local ordinance, exactly as $c$ is: black at our layer, and not necessarily black at the next.

That has a consequence I did not expect when I began this section. What falls in does not reach a bottom and does not vanish. It goes down a layer. Whatever the standard picture means when it worries that information is destroyed at a singularity, the tower's version is milder and stranger: nothing is destroyed. It is demoted --- kept, in a currency our instruments do not read.

Now the bill, because a section that only collects is a section this book does not print.

The first item is already in the ledger. It is the sixth wound, and this section is where I found it: if a body dense enough to trap light also stops the rain, its push follows its silhouette instead of its mass, and the sky says otherwise. I will not run the argument again here. I record only that it was this ladder that turned it up, and that of everything in the chapter it is the piece I hold most loosely.

The second item looked like a debt and turned out, on inspection, to be a question about bookkeeping rather than about the world. Black holes are credited with an entropy, and it is famously proportional to the \emph{area} of the horizon rather than to the volume within. The entropy section of this chapter counted arrangements in volumes, layer by layer, so a change of shape at a horizon wants explaining.

But recall what that counting was: never a fact about the region itself, always a fact about what the layer above can tell apart. And now notice \emph{where} a horizon is. It is not where the body is --- the sixth wound had to establish that the body sits far inside it, on a strut of its own. A horizon is the surface at which our layer's light gives out: which is to say, the surface beyond which our layer can no longer tell one arrangement from another. So of course the count our layer can make is a count across that surface. An area law is what layer-indexed entropy \emph{must} give at a censored boundary. It is a measure of our blindness, drawn on the edge of the blindness, and it would be strange if it came out any other shape.

Which suggests the puzzle was never quite the puzzle it appeared to be --- and the reader has met its twin already, a few pages back. The Gibbs paradox looked for a century and a half like a fact about gases and turned out to be a fact about the describer: the entropy of mixing appears or vanishes according to whether the experimenter can tell the gases apart. The area law has the same shape. It looks like a fact about what is inside a black hole, and on this book's reading it is a fact about where our light stops --- which is not where anything material is. The tower does not have to explain why the interior's states should be counted on a surface. It says they are not being counted at all; what is being counted is the boundary of our ignorance, and that boundary happens to be a sphere.

And the tower need not choose between the two numbers, because it has room for both. The area law is our layer's entropy of that region: what we, from outside, can distinguish. Whatever is down there has its own count at its own layer, and there is no reason on earth for the two to agree --- they are entropies of different descriptions, which is precisely what the coarse-graining ladder said all entropies are. The tower predicts that they differ, and predicts which is larger. I should say plainly that not everyone reads black-hole entropy this way; there are serious programmes that read it as a census of real interior states, and if that reading is right the tower owes an account of the interior and not merely of the view. But on the reading this book has used throughout, the area law is not a debt at all. It is the definition, doing what it was built to do.

One consequence, and I record it because it retires an idea of my own. I had thought the area law might be explained mechanically: the crowd itself stopped at the body's surface, the shell being all the outside can ever touch. That idea can be set down now. The area law belongs to the horizon, and the horizon is not the surface of anything material; a shell would have been answering a question that was never asked of it. What remains of the thought is not an asset but a hazard --- it is exactly the opacity that the sixth wound says would break the gravity.

So the section closes as the good ones do, with a gain and a bill on the same page. The gain is real: the tower's least testable claim --- structure below structure, without a bottom in sight --- turns out to be something the sky demonstrates twice over in objects we have photographed, and the one place where the standard picture openly admits an absurdity is the very place the tower does not need one. The bill is real too: the same collapsed bodies that make the case are where the mechanism is likeliest to fail. That is the trade this chapter keeps making, and I would rather own both halves than either.

